\section{introduction}
Over the past decade, artificial intelligence (AI) and machine learning (ML) have fundamentally reshaped data science, driving advances across domains as diverse as machine translation~\citep{isikscaling}, recommendation systems~\citep{yuan2025contextgnn}, social simulation ~\citep{yang2025twinmarket}, and medical diagnostics~\citep{sepehrimediconfusion}.
The growing availability of large-scale datasets~\citep{ghorbaniscaling}, coupled with rapid algorithmic progress, has enabled models that deliver increasingly accurate and adaptive outcomes.
However, as data becomes more heterogeneous and high-dimensional, so does the demand for experienced data scientists who can craft appropriate models, interpret nuanced patterns, and iterate toward optimal solutions.
% Over the past decade, artificial intelligence (AI) and machine learning (ML) have transformed the data science landscape, unlocking new possibilities for tackling problems across industries. From personalized e-commerce recommendations~\citep{ko2022survey, isinkaye2015recommendation} to AI-assisted healthcare diagnostics~\citep{khanna2022economics,shaheen2021applications}, organizations have leveraged vast datasets and ever-improving algorithms to deliver outcomes. Yet as data complexity grows, so does the demand for experienced data scientists who can craft appropriate models, interpret nuanced patterns, and iterate toward optimal solutions. 

% Crowdsourcing platforms like Kaggle\footnote{https://www.kaggle.com/} have emerged as a partial solution to this expertise bottleneck, bringing together thousands of data scientists to tackle challenging problems competitively. These platforms showcase both the immense value of human expertise and its inherent limitations: even elite competitors require substantial time in manual experimentation, feature engineering, and model tuning. Despite the collective intelligence of global participants, the iterative nature of data science workflows remains fundamentally constrained by human bandwidth and the trial-and-error process.

Crowdsourcing platforms like Kaggle\footnote{https://www.kaggle.com/} partially mitigate the expertise bottleneck by mobilizing thousands of data scientists, yet they also expose the limits of human-driven workflows: even top teams spend considerable time on trial-and-error experimentation, feature crafting, and hyperparameter tuning, making progress labor-intensive and slow.

% Large language models (LLMs) offer an opportunity to mitigate these limitations. With their demonstrated capabilities in code generation and reasoning~\citep{achiam2023gpt,team2023gemini,liu2024deepseek}, LLMs offer the potential to automate the exploration of solution spaces. %with substantially improved efficiency. 
Large language models (LLMs) offer an opportunity to mitigate these limitations. Their demonstrated strengths in code generation and reasoning\citep{achiam2023gpt,team2023gemini,liu2024deepseek} suggest they can automate exploration of large design spaces.
Machine Learning Engineering (MLE) serves as an ideal testbed for this potential, as success critically depends on efficiently navigating vast configuration spaces to identify optimal solutions. However, a significant gap persists between promise and practice.
Benchmarks like MLE-bench~\citep{chan2024mle-bench}, built on real Kaggle competitions, report that state-of-the-art LLMs reach only a small fraction of human expert performance, whereas well-designed agents can perform much better, highlighting the potential of improved agent design to accelerate MLE exploration.

Inspired by the common workflow of data scientist, we introduce R\&D-Agent, a comprehensive, decoupled, and extensible framework that formalizes the MLE process. Mirroring how data scientists work in practice, R\&D-Agent separates: (i) a \textbf{\uline{R}esearch phase} focused on idea generation and search-covering \textit{planning} (e.g., dynamically adjusting guidelines over time), \textit{exploration path structures} (e.g., tree-based search vs. chain-based search), \textit{memory context} (organizing and retrieving prior solutions and knowledge), and \textit{reasoning pipelines} (for hypothesis formation and refinement); and (ii) a \textbf{\uline{D}evelopment phase} focused on implementation and feedback-covering \textit{coding workflows} (steps turning ideas into runnable code efficiently) and \textit{evaluation strategy} (obtaining reliable and robust data-driven feedback). For each aspect, a simple and LLM-first baseline strategy is established as a clear starting point.

Positioning existing systems within this decomposition clarifies their coverage and gaps. As shown in Table~\ref{tab:agent_comparison}, previous work typically optimizes only a narrow slice of the workflow. For instance, AIDE~\citep{aide2025} and ML-Master~\citep{liu2025ml} primarily target the \textit{exploration path structure} via tree-based search, while MLE-STAR~\citep{nam2025mle} explores deeply along a single chain. KompeteAI~\citep{kulibaba2025kompeteai} emphasizes the \textit{coding workflow} with faster debugging. \textbf{All the existing methods can be summarized as a partial optimization from our framework's simple baseline.}


\begin{table}[h!]
\centering
\caption{Comparison of R\&D-Agent with existing methods, grouped by \textbf{Research} and \textbf{Development} phase design aspects.
``/'' denotes not covered or a simple, LLM-first baseline strategy.
  Research phase includes: (i) \textit{Planning}, (ii) \textit{Exploration Path Structuring}, (iii) \textit{Memory Context}, (iv) \textit{Reasoning Pipeline}.  
  Development phase includes: (v) \textit{Coding Workflow}, (vi) \textit{Evaluation Strategy}.  
  \faCog\ indicates that the framework provides multiple options followed by recommended best practices. 
  Details of R\&D-Agent's design are provided in Sec.~\ref{sec:method-design}.
}
\label{tab:agent_comparison}
\renewcommand{\arraystretch}{1.15}
\resizebox{\linewidth}{!}{%
\begin{tabular}{cl|cccc|cc}
\toprule
\multirow{3}{*}{\textbf{Phase}} &
\multirow{3}{*}{\textbf{Design Aspect}} &
\multicolumn{4}{c|}{\textit{Agent Designs}} &
\multicolumn{2}{c}{\textit{Frameworks}} \\
\cmidrule(lr){3-8}
& & AIDE & ML-Master & KompeteAI & MLE-STAR & AIRA & \textbf{R\&D-Agent} \\
& & \citep{aide2025} & \citep{liu2025ml} & \citep{kulibaba2025kompeteai} & \citep{nam2025mle} & \citep{toledo2025ai} & \textbf{(Ours)} \\
\midrule
\multirow{4}{*}{\textbf{Research}} 
& \textit{Planning}           & /               & /              & /                & /         & /                        & \faCog\ Dynamic \\
& \textit{Path Structuring}   & Tree (Greedy)   & Tree + MCTS    & Tree + Merging   & Chain     & \faCog\ Tree + MCTS & \faCog\ Adaptive \\
& \textit{Memory Context}     & /               & Sibling        & Sibling          & Sibling   & \faCog\ Sibling     & \faCog\ Collab. Comm.\\
& \textit{Reasoning Pipeline} & One step        & One step       & One step         & One step  & One step                 & \faCog\ Scientific multi-step
\\
\midrule
\multirow{2}{*}{\textbf{Development}} 
& \textit{Coding Workflow}    & Node Debug      & Node Debug     &Debug       & Debug  & /                       & \faCog\  Eff. \& Iter. Debug \\
& \textit{Evaluation Strategy}& /       & /        & /        & /          & /             & \faCog\ Aggregated \\
\bottomrule
\end{tabular}%
}
\end{table}

Although these systems report promising gains on their chosen components, they often (i) cover only a subset of the full MLE workflow, (ii) entangle multiple steps into monolithic pipelines rather than cleanly separating them into specialized components or agents, and (iii) leave many alternative designs within each phase unexplored. As a result, insights and improvements are difficult to generalize or reuse across tasks and domains. While \citet{toledo2025ai} recognize the need for a flexible framework, their design offers limited agent-level modularity and restricts exploration of the broader design space.

Our framework directly addresses these issues. R\&D-Agent (i) is comprehensive, enabling systematic exploration across the entire MLE workflow; (ii) is well-decoupled, allowing each phase to be implemented, replaced, and improved by dedicated components or expert agents; and (iii) is highly extensible, supporting plug-and-play alternatives that unify and generalize prior systems while facilitating the discovery of new agent configurations. In short, R\&D-Agent turns agent design for MLE from ad-hoc craftsmanship into a principled, testable process. 

Guided by human expertise, we instantiate R\&D-Agent and conduct systematic ablations across each phases to isolate the contribution of each component. Composing the best-performing choices shown in Table~\ref{tab:agent_comparison} within the framework yields a well-configured agent that delivers significant gains on MLE-bench~\citep{chan2024mle-bench}, achieving new state-of-the-art results.

% Many previous works have recognized this potential and proposed specific designs of the agent system to boost efficiency.
% As shown in Table~\ref{tab:agent_comparison}, this has led to diverse designs.
% Some methods focus on optimizing the \textit{exploration path structure} to reach good solutions more easily. For example, AIDE~\citep{aide2025} and ML-Master~\citep{liu2025ml} use tree-based structures, while MLE-STAR~\citep{nam2025mle} follows a chain-based approach to explore deeply along a single path.  
% Other methods, such as KompeteAI~\citep{kulibaba2025kompeteai}, emphasize the importance of development efficiency and introduce a novel \textit{coding workflow} that uses faster debugging strategies.

% These diverse methods optimize some key parts of the MLE workflow and achieved promising results,  
% they still have notable shortcomings:  
% 1) they cover only a subset of the complete MLE workflow;  
% 2) they combine several key steps into a single monolithic pipeline instead of clearly separating them into dedicated components or specialized agents (MLE is a complex task that benefits from the collaborative intelligence of multiple expert agents);
% 3) besides, for each component, there are alternative designs that have not yet been explored.
% These gaps leave substantial room for further investigation.  
% To systematically explore and realize this potential, a flexible MLE-agent framework becomes crucial for efficiently exploring the optimal design of agent systems. \cite{toledo2025ai} recognized this potential and proposed a framework providing significant insights, but offering limited agent design flexibility.


% To address these challenges, we further analyze the scenarios and introduce a more comprehensive, decoupled, and extensible framework.
% As a scientific discipline, MLE requires sustained and focused research and development to achieve strong results.  
% Inspired by the typical workflow of data scientists, we divide the process for a capable MLE Agent into two main aspects:
% (i) thinking and proposing ideas, which we define as the \textbf{research phase}. This includes \textit{planning} (e.g., dynamically adjusting guidelines over time), using efficient \textit{exploration path structures}, organizing \textit{memory context} to make better use of past experiences and knowledge, and applying various \textit{reasoning pipelines} for idea generation; and  
% (ii) turning these ideas into runnable code and getting evaluation feedback from data, which we define as the \textbf{development phase}. This involves implementing ideas efficiently through well-designed \textit{coding workflows} and ensuring that the \textit{evaluation strategy} provides reliable and robust feedback from the data.

% \todo{the definition of scenarios is still not easy to understand}

% We argue that breaking this architectural coupling through explicit phase separation is essential for breakthrough performance. The nature of MLE tasks inherently requires agents to navigate a complex design space composed of several key aspects: strategic components such as \textit{planning}, which allocates exploration effort over time; \textit{exploration path structuring}, which organizes how solutions are explored; \textit{memory context}, for managing accumulated knowledge; and a \textit{reasoning pipeline} to generate new hypotheses. Equally crucial are the tactical components: a \textit{coding workflow} to translate ideas into robust code, and an \textit{evaluation strategy} to ensure reliable performance assessment.

% These complementary design aspects naturally group into two distinct phases. The strategic components collectively form the \textbf{Research Phase}, which addresses the overarching question of ``what to try''. The tactical components constitute the \textbf{Development Phase}, which tackles the question of ``how to implement''. To systematically explore how to best implement this principle, a flexible MLE-agent framework becomes crucial. While \cite{toledo2025ai} recognized this need and provided valuable insights through their framework, it still offers limited design flexibility.

% The primary challenge for current MLE agents is their architectural rigidity. Unlike human experts who fluidly separate strategic exploration from tactical implementation, existing agents are trapped in monolithic pipelines that conflate these distinct cognitive modes. 
% Both paradigms force agents to simultaneously handle strategy and implementation, 
% creating a cognitive burden that fundamentally limits exploration efficiency and solution quality.

% We argue that breaking this architectural coupling through explicit phase separation is essential for breakthrough performance. The nature of MLE tasks inherently requires two distinct cognitive modes: (i) a \textbf{Research Phase} for addressing "what to try" through strategic planning, exploration structuring, and hypothesis generation; and (ii) a \textbf{Development Phase} for tackling "how to implement" through efficient coding, debugging, and evaluation. Successful agents must recognize and operationalize this natural separation. To systematically explore how to best implement this principle, a flexible MLE-agent framework becomes crucial. While \cite{toledo2025ai} recognized this need and provided valuable insights through their framework, it still offers limited design flexibility.


% To fully explore the potential of LLM-based MLE agent design,  
% To address this challenge, we propose a general and highly extensible MLE-Agent framework called R\&D-Agent.  
% Within this framework, we test a wide range of methods and provide valuable insights into agent design across diverse MLE tasks, ultimately identifying a design that achieves state-of-the-art performance.

% Based on this intuition, we present R\&D-Agent, a framework for MLE agents that unifies and generalizes all known previous MLE agent systems listed in Table~\ref{tab:agent_comparison}.  
% 1) Our framework is comprehensive, enabling systematic exploration of the MLE agent design space;  
% 2) it is well-decoupled, allowing for the dedicated implementation of key components and enabling deeper intelligence;  
% 3) it is also highly extensible, supporting the discovery of novel agent configurations and expanding the boundaries of agent research.  
% This design transforms agent development from an art into a science.
% % Unlike existing fixed architectures, our framework offers composable modules that make it easy to explore different MLE agent designs systematically.
% Extensive experiments on MLE-bench show that our comprehensive, decoupled, and extensible framework delivers significant performance gains, with our best configuration achieving new state-of-the-art results.

% This modularity transforms agent design from art into science, allowing researchers to: 1) reproduce and compare existing approaches within a unified framework, 2) conduct controlled experiments to identify key performance factors, and 3) discover novel agent configurations that surpass current methods. 

In summary, our primary contributions are:
\begin{itemize}[leftmargin=4ex]
  \item We introduce R\&D-Agent, a comprehensive, decoupled, and extensible autonomous agent framework that formalizes the MLE process by separating a research phase (planning, exploration path structure, memory context, and reasoning pipelines) from a development phase (coding workflows and evaluation strategy). 
  \item With R\&D-Agent enabling plug-and-play alternatives, we conduct systematic ablations across both phases to isolate the contribution of each component and derive insights into the factors that most affect agent performance on data science.
  \item Guided by human expertise, an R\&D-Agent configuration that composes the best-performing choices within the framework delivers significant gains on MLE-bench~\citep{chan2024mle-bench}, achieving new state-of-the-art results and demonstrating the efficiency and effectiveness of our approach.
  % \item We identify exploration efficiency as the main bottleneck for MLE agents and propose comprehensive and dedicated modules as a core design principle. By separating the strategic research phase from the tactical development phase, we present a novel framework named R\&D-Agent, which greatly improves the efficiency of discovering and validating agent-system designs.
  % \item With R\&D-Agent, an autonomous and configurable framework, we gain valuable insights into the key components that impact agent performance.
  % \item An R\&D-Agent-based design, inspired by the data science workflow of human experts, achieves state-of-the-art performance on MLE-bench under stricter time constraints than prior work, clearly demonstrating the superior efficiency and effectiveness of our framework.
\end{itemize}

% However, existing designs still overlook several key challenges and cover only a limited range of agent system architectures.  
%
% (i) First, while chain-structured agents could be more efficient than tree-based designs, they are susceptible to suboptimal solutions due to hallucinations or the need for very long memory over many iterations. If the reasoning capabilities of the underlying LLM are sufficiently strong, however, \textbf{chain-based structures} can outperform tree-based approaches in terms of efficiency. On the other hand, \textbf{Tree-based structures}, such as those employing MCTS, typically require more iterations, and the ideas generated at each step may not fully leverage the reasoning capabilities of the LLM, often resulting in suboptimal or mediocre solutions.
%
% (ii) Second, the execution time of code constitutes a significant portion of overall agent runtime. Few agents explicitly consider the efficiency of code execution, yet it is a critical factor affecting the overall performance of the system.
%
% (iii) Third, the scalability of exploration strategies remains very limited in most agents. In particular, chain-based structures are difficult to directly adapt into tree-based searches, and tree-based approaches cannot readily inherit the advantages of chain-based reasoning. Moreover, combining the two approaches in a complementary manner could potentially enable more efficient exploration paths, yet existing agents rarely implement such hybrid strategies.


% A key feature of R\&D-Agent v2 is its \textbf{Modular Directed Execution Graph} (MDEG), which serves as the backbone of the agent loop. This DAG-based design provides a highly flexible and extensible structure that can naturally encompass both chain-based and tree-based agent architectures. Specifically, a chain structure can be represented as a simple linear path within the DAG, while a tree structure corresponds to a branching configuration where modules have multiple successors. This unified representation allows R\&D-Agent to combine the advantages of both paradigms, supporting fast idea generation through chain-like paths and broad exploration through tree-like branches, all while maintaining a clear execution order without cycles.  

% To highlight the distinctions between our framework and prior approaches, we summarize them in the table below.
% \begin{table}[h!]
% \centering
% \caption{Comparison of R\&D-Agent v2 with existing methods.}
% \label{tab:agent_comparison}
% \resizebox{\linewidth}{!}{%
% \begin{tabular}{lcccc} % l: left, c: center
% \hline
% \textbf{Method} & \textbf{Structure} & \textbf{Reasoning Util.} & \textbf{Code Eff.} & \textbf{Scal.} \\
% \hline
% AIDE \citep{aide2025} & Tree (Greedy) & Medium & Low & Limited \\
% ML-Master \citep{liu2025ml} & Tree + MCTS & High & Medium & Limited \\
% KompeteAI \citep{kulibaba2025kompeteai} & Tree + Merging & High & Medium & Limited \\
% MLE-STAR \citep{nam2025mle} & Chain & High & Medium & Medium \\
% \textbf{R\&D-Agent v2 (Ours)} & MDEG (Flex./Hybrid) & Very High & High & High \\
% \hline
% \end{tabular}%
% }
% \end{table}

% \begin{table}[h!]
% \centering
% \caption{Comparison of R\&D-Agent with existing methods, grouped by \textbf{Research} and \textbf{Development} phase design aspects. ``/'' denotes not covered or naive.  
%   Research phase includes: (i) \textbf{Planning}, (ii) \textbf{Exploration Path Structuring}, (iii) \textbf{Memory Context}, (iv) \textbf{Reasoning Pipeline}.  
%   Development phase includes: (v) \textbf{Coding Workflow}, (vi) \textbf{Evaluation Strategy}.  
%   Methods are categorized as \textit{Concrete Designs} or \textit{Frameworks}.  
%   \faCog\ indicates that the framework provides multiple options followed by recommended best practices.  
% }
% \label{tab:agent_comparison}
% \resizebox{\linewidth}{!}{%
% \begin{tabular}{lcccccc} % l: left, c: center
% \hline
% \textbf{Method} & \multicolumn{4}{c|}{\textbf{Research Phase}} & \multicolumn{2}{c}{\textbf{Development Phase}} \\
% \cline{2-5} \cline{6-7}
%  & \textbf{Plan} & \textbf{Path Struct.} & \textbf{Mem. Context} & \textbf{Reasoning Pipe.} & \textbf{Coding} & \textbf{Eval.} \\ 
% \hline
% \multicolumn{7}{l}{\textit{Agent Designs}} \\
% \hline
% AIDE \citep{aide2025}  & /  & Tree (Greedy) & / & One step reasoning & Node Debug & /   \\
% ML-Master \citep{liu2025ml} & /  & Tree + MCTS & Sibling & One step reasoning & Node Debug & /   \\
% KompeteAI \citep{kulibaba2025kompeteai} & / & Tree + Merging &   Sibling & One step reasoning & Eff. Debug & /  \\
% MLE-STAR \citep{nam2025mle} & / & Chain & Sibling & One step reasoning& Eff. Debug & / \\
% \hline
% \multicolumn{7}{l}{\textit{Agent Frameworks}} \\
% \hline
% AIRA \citep{toledo2025ai} & / & \faCog\Tree + MCTS & \faCog\ Sibling &  One step reasoning & / & /  \\
% \textbf{R\&D-Agent (Ours)} & \faCog\ Dynamic  &  \faCog\ Adaptive Explore \& Exploit & \faCog\ Collab. Comm. & \faCog\ Scientific  & \faCog\ Eff. Debug & \faCog\ Multi.\\
% \hline
% \end{tabular}%
% }
% \end{table}






% \todo{the table is not accurate enough}
% \todo{the words in the table is not well explained}

% Within each loop of the DAG, we define several fundamental modules: \textit{Hypothesis Gen}, \textit{Task Gen}, \textit{Coder}, \textit{Runner}, \textit{Feedback}, and \textit{Record}. Among these, \textit{Hypothesis Gen} is the core module responsible for generating new ideas. \textit{Task Gen}, \textit{Coder}, and \textit{Runner} work together to efficiently implement and execute code, enhancing overall computational efficiency. Finally, \textit{Feedback} and \textit{Record} act as memory units, capturing outcomes and informing subsequent iterations. This modular DAG design not only provides high flexibility and scalability but also enables efficient exploration and execution within the agent system.





\begin{comment}
Crucially, data science projects often hinge on iterative insight -- the repeated cycle of exploring, testing, and refining approaches as new knowledge emerges. Even when two projects share the same overarching objective, variations in data distribution, sample size, or domain constraints can demand radically different paths towards optimal solutions. Expert data scientists use feedback from each experimental iteration -- whether on feature importance, model fit, or resource limitations -- to systematically adapt their methods. Because each iteration can be expensive and time-consuming, these projects place a premium on well-conceived exploration strategies rather than relying on random or brute-force attempts.


In light of these challenges, we propose that a successful machine learning engineering agent must actively learn and adapt through iterative exploration. It should synthesize domain insights, generate deeper hypotheses, and refine its approach based on partial findings rather than rely on a single solution path. With these goals in mind, this paper introduces R\&D-Agent, a novel framework designed to drive more effective exploration in data-driven projects across a wide array of industries. By boosting both the efficiency of research exploration and the precision of development, R\&D-Agent aspires to narrow the gap between automated intelligence and expert-level data science, accelerating innovation where it is needed most.

In particular, the design of the R\&D-Agent is guided by two key principles:
\begin{enumerate}
\item \emph{Dedicated R\&D Role}: the framework employs two specialized agents -- the ``Researcher'' and the ``Developer'' -- which correspond to the two types of feedback provided in each exploration step: solution performance and execution error information. The Researcher agent processes performance feedback to drive efficient exploration (i.e., idea generation), while the Developer agent leverages execution logs to iteratively refine solution implementation (i.e., code generation). Dedicated agents, designed for specific tasks and equipped with corresponding feedback, facilitate more efficient exploration.
%\item \emph{Decomposed Workflow}: the framework decomposes a data science solution into distinct, decoupled components -- such as feature generation, model training, and ensemble -- organized into a pipeline, with each component linked to targeted improvement ideas. To enhance exploration efficiency, we propose applying multiple independent improvement ideas concurrently in subsequent iterations. This design -- combining solution decomposition with idea alignment -- ensures that enhancements to one component can be implemented simultaneously with improvements to others, without interference, ultimately yielding a more significant overall advancement in a single exploration step. Moreover, ideas aligned with specific components have focused objectives and can more effectively leverage domain knowledge from human experts, ultimately leading to more efficient exploration.
\item \emph{Mutually-Enhanced Multiple Traces}: the framework supports multiple exploration traces running in parallel while facilitating advanced mutual enhancements. For instance, you can launch a new exploration trace from an existing checkpoint or merge checkpoints from different traces to generate a more powerful composite solution. These capabilities not only enable the parallel exploration of a single data science task but also enhance the performance of individual exploration traces, thereby significantly boosting overall exploration efficiency.
\end{enumerate}

Notably, the framework's design offers significant flexibility for integrating diverse research contributions. This capability not only boosts the efficiency of research exploration but also enhances development precision -- all through simple API calls provided by our R\&D-Agent framework. Ultimately, this synergy results in a more robust data science automation system.
\end{comment}

